\subsection{Scenarios}
Missions are created by a seeded generator with the parameters in Table~\ref{tab:params}. The map is a square with a
side of 10\,km and the depot at its centre, and half of the tasks are clustered around three hotspots. Half of the tasks
are known at the start, and the others arrive in batches. Three fleet scales with 5, 10 and 20 UAVs and 25, 50 and 140
tasks are combined with three event profiles. The arrivals profile contains six arrival batches, the failures profile
three batches and two UAV failures, and the deadlines profile three batches and four tightened deadlines. A fourth scale
with 40 UAVs and 280 tasks is used for scalability only.

\begin{table}[t]
\caption{Scenario generator parameters (distances in m, times in s, energy in metres of flight).}\label{tab:params}
\centering\small
\setlength{\tabcolsep}{4pt}
\begin{tabular}{ll@{\qquad}ll}
\toprule
Map side / horizon & 10\,000 / 1\,800 & Service time & $U[30,90]$ \\
UAV speed & $U[10,16]$\,m/s & Time-window slack & $U[180,600]$ \\
Energy capacity & $U[18,28]$\,km & Release delay & $U[0,120]$ \\
Energy reserve & 10\% of capacity & Rewards 1, 2, 3, 5 & probabilities .4, .3, .2, .1 \\
Hover cost & 6 per second & Hotspots & 3, std.\ 600, 50\% of tasks \\
Tasks known at start & 50\% & Scales (UAVs/tasks) & 5/25, 10/50, 20/140, 40/280 \\
\bottomrule
\end{tabular}
\end{table}

\subsection{Static snapshots with exact execution}
The stopping rules are first compared without execution uncertainty. Each mission is simulated once with a
deterministic reference planner, and every replanning call after the start is stored with its previous plan. All
methods are evaluated on these snapshots, and regret is measured against the best known objective of each snapshot,
which includes a 3\,s ALNS run and, for the small and medium scales, the OR-Tools routing solver~\cite{ortools} with
guided local search~\cite{voudouris1999gls}. Because all ALNS variants follow the same search trajectory, whole families
of each rule are evaluated exactly offline, with 12 fixed budgets, 10 stagnation patience values and 9 compute prices,
and each baseline is compared with AnytimeSwarm at the same mean computation. This benchmark uses its own planned
test seeds 2000--2029 (\val{NsnapTest} snapshots of 270 missions), and it was reproduced on the final software platform
before the execution-aware study (Section~\ref{sec:repro}).

